%
% teil3.tex -- Beispiel-File für Teil 3
%
% (c) 2020 Prof Dr Andreas Müller, Hochschule Rapperswil
%
% !TEX root = ../../buch.tex
% !TEX encoding = UTF-8
%
\section{Fazit
\label{Fazit}}

Die Aussage, dass die Summe aller natürlichen Zahlen -1/12 ist, macht
auf den Ersten Blick keinen Sinn, da die Reihe offensichtlich
divergiert. Dennoch ist der Wert -1/12 im Rahmen der analytischen
Fortsetzung der Zeta-Funktion mathematisch belegt und sinnvoll
definiert.

Bei dem Ergebnis, das man von der analytischen Fortsetzung der
Riemannschen Zeta-Funktion erhält, handelt es sich jedoch nicht um eine
Summe im herkömmlichen Sinn. Es ist eher als erweiterte Interpretation
zu betrachten, die in mathematischen Kontexten nützlich ist.


